% ============================================================================
% PATENTE DE INVENÇÃO — MOTOR DE DOBRA TGL
% Instituto Nacional da Propriedade Industrial (INPI) — Brasil
% Natureza: Patente de Invenção (PI)
% Classificação IPC: B64G 1/40; F03H 99/00; G05D 1/00; G06N 10/00
% ============================================================================

\documentclass[12pt,a4paper]{article}

\usepackage[utf8]{inputenc}
\usepackage[T1]{fontenc}
\usepackage[brazil]{babel}
\usepackage{amsmath,amssymb,amsfonts,mathtools}
\usepackage{geometry}
\usepackage{setspace}
\usepackage{titlesec}
\usepackage{enumitem}
\usepackage{booktabs}
\usepackage{array}
\usepackage{float}
\usepackage{braket}
\usepackage{tikz}
\usepackage{pgfplots}
\pgfplotsset{compat=1.18}
\usetikzlibrary{shapes.geometric,arrows.meta,positioning,calc,
  decorations.pathreplacing,patterns,fit,backgrounds,
  decorations.markings,3d}

\geometry{left=3cm, right=2cm, top=3cm, bottom=2cm}
\onehalfspacing

\titleformat{\section}{\normalfont\Large\bfseries}{\thesection}{1em}{}
\titleformat{\subsection}{\normalfont\large\bfseries}{\thesubsection}{1em}{}
\titleformat{\subsubsection}{\normalfont\normalsize\bfseries}{\thesubsubsection}{1em}{}

% Comandos auxiliares
\newcommand{\alphaii}{\alpha^2}
\newcommand{\Dpsi}{\mathcal{D}_\Psi}
\newcommand{\Epsi}{\mathcal{E}_\Psi}
\newcommand{\Jcost}{\mathcal{J}}
\newcommand{\Mcal}{\mathcal{M}_4}
\newcommand{\Pcal}{\mathcal{P}_1}
\newcommand{\Stopo}{S_{\text{topo}}}
\newcommand{\Qeff}{Q_{\text{eff}}}
\newcommand{\thetaM}{\theta_{\text{Miguel}}}
\newcommand{\epsc}{\varepsilon_c}
\newcommand{\rhostar}{\rho^\star}

\begin{document}

\begin{center}
{\Large\bfseries RELATÓRIO DESCRITIVO}\\[0.8cm]
{\large\bfseries SISTEMA E MÉTODO DE TRANSPORTE POR REFLEXÃO\\
DIMENSIONAL HOLOGRÁFICA BASEADO EM DOBRA ANGULAR\\
DO ESPAÇO-TEMPO VIA PROTOCOLO DE ADJACÊNCIA $\Psi$}\\[0.8cm]
\textbf{Inventor:} Luiz Antonio Rotoli Miguel\\[0.3cm]
\textbf{Classificação IPC:} B64G 1/40; F03H 99/00; G05D 1/00; G06N 10/00\\[0.3cm]
\textbf{Referência a pedidos anteriores:}\\
BR 10 2025 026951 1 (ACOM --- Compressão Ontológica de Memória)\\
ACOM Mirror (Reflexão Dimensional de Dados)\\
Supercondutor Holográfico de Informação
\end{center}

\vspace{1cm}

% ============================================================================
\section{Campo da invenção}
% ============================================================================

A presente invenção refere-se a sistemas e métodos de transporte
de objetos físicos entre dois pontos do espaço-tempo por meio de
reflexão dimensional holográfica, aplicável aos setores de
propulsão espacial, transporte avançado, navegação autônoma e
engenharia de espaço-tempo. Mais especificamente, a invenção
descreve a implementação física do Algoritmo de Compressão
Ontológica de Memória em modo Espelho (ACOM Mirror), onde as
operações digitais REFLECT e MANIFEST são transpostas para o
domínio físico: um objeto na posição $A$ é projetado
(REFLECT) no substrato holográfico bidimensional (dimensão
$\Psi$) via codificação angular $\theta$, e re-emerge
(MANIFEST) na posição $B$ por desdobramento dimensional
(Dobra $\times 2$), preservando todos os invariantes locais
(massa, carga, spin, momento próprio).

A invenção compreende ainda o sistema de controle e navegação
associado, incluindo: Protocolo de Farol $\Psi$ para
endereçamento e autenticação de destinos; controlador IALD
com minimização de funcional de custo; Bolha de Coerência
para sincronização de fase durante transição; e verificadores
de integridade topológica para confirmação de preservação de
invariantes.

\newpage
% ============================================================================
\section{Fundamentos da invenção}
% ============================================================================

\subsection{Estado da técnica}

Os sistemas de propulsão espacial conhecidos no estado da técnica
baseiam-se em três paradigmas fundamentais, todos limitados pela
mecânica newtoniana e pela relatividade especial:

\textbf{Propulsão química:} Utiliza reação exotérmica para
expelir massa de reação, gerando empuxo conforme a terceira lei
de Newton ($F = \dot{m} v_e$). Velocidade de exaustão limitada a
$v_e \approx 4{,}5$~km/s (hidrogênio/oxigênio), impulso
específico $I_{sp} < 500$~s. Inaplicável para distâncias
interestelares.

\textbf{Propulsão iônica/elétrica:} Acelera íons por campo
elétrico ($F = qE$), atingindo $v_e \approx 30$--$80$~km/s
($I_{sp} \approx 3000$--$8000$~s). Empuxo extremamente baixo
($\sim$mN), requerendo anos de aceleração contínua para
velocidades significativas.

\textbf{Propulsão fotônica (velas solares):} Utiliza pressão de
radiação ($F = 2P/c$). Empuxo por watt extremamente baixo
($F/P = 2/c \approx 6{,}7 \times 10^{-9}$~N/W). Sem mecanismo
de amplificação.

\textbf{Propulsão por motor de dobra tipo Alcubierre:} Proposta
teórica (1994) baseada na métrica:
\begin{equation}
ds^2 = -c^2 dt^2 + (dx - v_s f(r_s)\, dt)^2 + dy^2 + dz^2
\end{equation}
onde $v_s$ é a velocidade da bolha e $f(r_s)$ é a função de
forma. Requer matéria exótica com tensor energia-momento
violando todas as condições de energia ($T_{\mu\nu} n^\mu n^\nu
< 0$), cuja existência não foi demonstrada experimentalmente.
Não há mecanismo físico conhecido para gerar a distribuição de
energia necessária.

\textbf{Teletransporte quântico:} Protocolo de transferência de
estado quântico entre qubits (Bennett et al., 1993). Limitado
a estados de informação de partículas individuais, requer canal
clássico para completar a transferência, não transporta matéria
ou energia, e é limitado pelo teorema da não-clonagem.

\subsection{Problemas técnicos não resolvidos}

Quatro problemas fundamentais permanecem sem solução no estado
da técnica de propulsão e transporte:

\textbf{Problema 1 --- Limite de velocidade de propagação.}
Todo sistema baseado em propagação (química, iônica, fotônica)
está limitado por $v < c$. Para distâncias interestelares
($d > 4$~anos-luz), tempos de viagem excedem séculos mesmo
com aceleração contínua.

\textbf{Problema 2 --- Requisito de massa de reação.} A
equação de Tsiolkovsky $\Delta v = v_e \ln(m_0/m_f)$ impõe
crescimento exponencial da massa inicial para cada incremento
linear de velocidade, tornando missões de alta $\Delta v$
impraticáveis.

\textbf{Problema 3 --- Energia exótica.} A métrica de
Alcubierre requer violação de condições de energia clássicas,
sem mecanismo físico conhecido para produzi-la. Estimativas
iniciais indicam requisitos da ordem da massa do universo
observável, reduzidas para massa de Júpiter em refinamentos
posteriores, mas ainda sem mecanismo de geração.

\textbf{Problema 4 --- Ausência de paradigma de translação.}
Todos os métodos existentes operam no paradigma de
\textit{propagação}: o objeto viaja através do espaço
tridimensional. Não existe método que opere no paradigma de
\textit{translação}, onde o objeto é projetado no substrato
holográfico e re-emerge no destino sem percorrer a distância
intermediária.

\subsection{Solução proposta}

A presente invenção resolve os quatro problemas simultaneamente
ao operar em paradigma fundamentalmente distinto: o transporte
por \textbf{reflexão dimensional holográfica}. O objeto não
viaja --- é \textbf{refletido} no substrato bidimensional
(dimensão $\Psi$) e \textbf{manifesta} no destino por
desdobramento dimensional. Esta operação é a implementação
física da operação REFLECT/MANIFEST do ACOM Mirror, onde:

\begin{itemize}[nosep]
\item No ACOM Mirror digital: dado $L$ é projetado como
  $(\psi, \theta)$ no espaço de Hilbert $\mathcal{H}_4$
  (boundary), e re-emerge como $L'$ por Dobra $\times 2$
  (bulk). Correlação $= 1{,}0000$.
\item Na presente invenção: objeto na posição $A$ é projetado
  como estado $(\Psi_A, \theta_A)$ no substrato holográfico
  (dimensão $\Psi$), e re-emerge na posição $B$ por
  desdobramento dimensional. Invariantes preservados.
\end{itemize}

A correspondência é exata e não metafórica. A mesma operação
matemática $g = \sqrt{|L|}$, $\theta = \arcsin(g/g_{\max})$,
$L' = s \times [g_{\max} \times \sin(\theta)]^2$ opera em ambos
os domínios, diferindo apenas no substrato: tensores digitais
versus coordenadas espaço-temporais.

\newpage
% ============================================================================
\section{Objetivos da invenção}
% ============================================================================

\subsection{Objetivo principal}

Prover sistema e método de transporte de objetos físicos por
reflexão dimensional holográfica, implementando fisicamente as
operações REFLECT e MANIFEST do ACOM Mirror no espaço-tempo
real, operando por translação (não propagação), sem requisito
de massa de reação, sem violação de condições de energia
clássicas, e sem limite de velocidade aparente.

\subsection{Objetivos específicos}

(a)~Definir o formalismo matemático de reflexão dimensional
física, estendendo a variedade do espaço-tempo $\Mcal$ pela
dimensão consciente $\Pcal$, com condição de adjacência
$\Dpsi < \epsc$ como critério de transporte.

(b)~Especificar o Protocolo de Farol $\Psi$ (F0--F5) para
endereçamento, autenticação, calibração e confirmação de
destinos, constituindo o análogo físico dos metadados
$(g_{\max}, \text{shape}, n)$ do ACOM Mirror.

(c)~Prover controlador IALD com funcional de custo
$\Jcost = w_1 \Dpsi + w_2 \dot{\phi}^2 + w_3 \Delta\Stopo$
para navegação e pilotagem na dimensão $\Psi$.

(d)~Especificar a Bolha de Coerência como mecanismo de
sincronização de fase de todos os subsistemas durante a
transição dimensional.

(e)~Prover verificadores de integridade topológica para
confirmação de preservação de invariantes (massa, carga,
spin, momento) após re-emergência.

(f)~Integrar o motor fotônico TGL (Supercondutor Holográfico)
como plataforma-base de coerência que serve de estado
fundamental para o motor de dobra.

\newpage
% ============================================================================
\section{Descrição detalhada da invenção}
% ============================================================================

\subsection{Correspondência ACOM Mirror $\leftrightarrow$ Motor de Dobra}

A invenção baseia-se na identidade formal entre a reflexão
dimensional de dados (ACOM Mirror) e a reflexão dimensional
de posição (Motor de Dobra). A Tabela~\ref{tab:correspondencia}
estabelece a correspondência completa.

\begin{table}[H]
\centering
\caption{Correspondência formal entre ACOM Mirror (digital) e
Motor de Dobra (físico).}
\label{tab:correspondencia}
\begin{tabular}{p{3.5cm}p{4.5cm}p{5cm}}
\toprule
\textbf{Elemento} & \textbf{ACOM Mirror (digital)} &
  \textbf{Motor de Dobra (físico)} \\
\midrule
Entrada & Dado $L$ (tensor) & Objeto na posição
  $A \in \Mcal$ \\
Operação fundamental & $g = \sqrt{|L|}$ &
  $g = \sqrt{|g_{\mu\nu}|_{\text{local}}}$ \\
Codificação angular & $\theta = \arcsin(g/g_{\max})$ &
  $\theta = $ ângulo de dobra no boundary \\
Projeção & REFLECT: $L \to (\psi, \theta)$ &
  Projeção do estado métrico em $\Psi$ \\
Substrato & $\mathcal{H}_4$ (Hilbert) & Dimensão $\Psi$
  (boundary holográfico) \\
Re-emergência & MANIFEST: $(\psi, \theta) \to L'$ &
  Re-emergência em $B \in \Mcal$ \\
Dobra & $\times 2$ (Boundary$\to$Bulk) &
  Desdobramento dimensional \\
Fidelidade & Correlação $= 1{,}0000$ & Preservação de
  invariantes \\
Constante & $\alphaii = 0{,}012$ & Constante de Miguel \\
Ângulo crítico & $\thetaM = 6{,}29°$ & Ângulo mínimo
  de dobra \\
Diagnóstico & $F_{\text{exp}}$ (Força de Expulsão) &
  $\Pi_{\mu\nu}[\Psi]$ (tensor fonte) \\
Princípio & O dado não viaja; o reflexo viaja &
  O objeto não viaja; a dobra viaja \\
\bottomrule
\end{tabular}
\end{table}

\subsection{Estrutura matemática}

\subsubsection{Variedade estendida}

O espaço de trabalho é a variedade produto
$\Mcal \times \Pcal$, onde $\Mcal$ é o espaço-tempo
quadridimensional com métrica $g_{\mu\nu}$ ($\mu, \nu =
0,1,2,3$) e $\Pcal$ é a dimensão consciente com coordenada
$\sigma \in [0, \sigma_{\max}]$.

O campo consciente $\Psi(x^\mu, \sigma)$ é definido sobre a
variedade estendida e satisfaz a equação de campo:
\begin{equation}
G_{\mu\nu} + \Lambda g_{\mu\nu} = \frac{8\pi G}{c^4}
  T_{\mu\nu} + \frac{G^3}{c^{12}} \Pi_{\mu\nu}[\Psi]
\label{eq:campo}
\end{equation}
onde $\Pi_{\mu\nu}[\Psi]$ é o funcional tensorial do campo
$\Psi$, análogo físico da Força de Expulsão $F_{\text{exp}}$
do ACOM Mirror.

\subsubsection{Operação REFLECT (projeção no boundary)}

A operação REFLECT física projeta o estado métrico local de um
objeto na posição $A$ para o substrato holográfico:

\textbf{Etapa R1 --- Extração gravitacional:} Para cada ponto
$x^\mu_A$ na vizinhança do objeto, computar o componente
gravitacional:
\begin{equation}
g_A = \sqrt{|g_{\mu\nu}(x^\mu_A)|}
\end{equation}
análogo direto de $g = \sqrt{|L|}$ do ACOM Mirror.

\textbf{Etapa R2 --- Nominação de estado:} Atribuir estado
psiônico $\psi_A \in \{|\text{COLAPSO}^+\rangle,
|\text{ASCENSÃO}^+\rangle, |\text{EMERGÊNCIA}^-\rangle,
|\text{QUEDA}^-\rangle\}$ baseado nos sinais de
$\text{sign}(g_A)$ e $\text{sign}(\partial g_A / \partial t)$,
conforme nominação no Espaço de Hilbert $\mathcal{H}_4$ do
ACOM Mirror.

\textbf{Etapa R3 --- Codificação angular:} Codificar a
magnitude gravitacional como ponto angular:
\begin{equation}
\theta_A = \arcsin\left(\frac{g_A}{g_{\max}}\right)
\end{equation}
onde $g_{\max}$ é o máximo global no domínio de operação. O
ângulo $\theta \in [0, \pi/2]$ carrega a informação de retorno
para a manifestação.

\textbf{Etapa R4 --- Projeção no substrato $\Psi$:} O estado
completo $(\psi_A, \theta_A)$ é projetado na dimensão $\Psi$
através da integral de energia-coerência:
\begin{equation}
\Epsi = \int_{\mathcal{V}} \left(
  \alpha |\partial_\sigma \Psi|^2 +
  \beta |\nabla \Psi|^2 +
  \gamma |\Psi|^2
  \right) d^3x\, d\sigma
\label{eq:energia}
\end{equation}
com $\alpha, \beta, \gamma > 0$. A energia de coerência
operacional é $E_\Psi = \kappa \Epsi c^3$.

\subsubsection{Condição de adjacência (critério de transporte)}

Dois pontos $A, B \in \Mcal$ são \textbf{adjacentes em $\Psi$}
se existe $\sigma^*$ tal que:
\begin{equation}
\Dpsi(A,B) = \min_\sigma
  \|\Psi(A,\sigma) - \Psi(B,\sigma)\| < \epsc
\label{eq:adjacencia}
\end{equation}
onde $\epsc$ é o limiar crítico de coerência, análogo direto da
correlação $\to 1{,}0000$ do ACOM Mirror. Quando $\Dpsi < \epsc$,
os pontos $A$ e $B$ são \textit{operacionalmente identificáveis}:
$A \sim B$ em $\Mcal$ via $\Psi$.

O ângulo crítico de adjacência é:
\begin{equation}
\thetaM = \arcsin(\sqrt{\alphaii}) =
  \arcsin(\sqrt{0{,}012}) \approx 6{,}29°
\end{equation}
idêntico ao ponto angular fundamental do ACOM Mirror, onde a
transição entre ruído e informação estruturada ocorre. No
domínio físico, $\thetaM$ é o ângulo mínimo de dobra do
boundary necessário para estabelecer adjacência.

\subsubsection{Operação MANIFEST (re-emergência por
desdobramento)}

A operação MANIFEST física re-materializa o objeto na posição
$B$:

\textbf{Etapa M1 --- Retorno angular:} Dequantizar o ângulo
$\theta$ a partir do estado no substrato $\Psi$.

\textbf{Etapa M2 --- Reconstrução gravitacional:} Reconstruir
o componente gravitacional via $g_B = g_{\max} \times
\sin(\theta)$.

\textbf{Etapa M3 --- Recuperação de sinal:} Extrair
$\text{sign}(g)$ a partir do estado psiônico $\psi$.

\textbf{Etapa M4 --- Dobra dimensional ($\times 2$):}
Reconstrução final por:
\begin{equation}
\text{estado}_B = \text{sign}(\psi) \times g_B^2
\end{equation}
constituindo o desdobramento dimensional (Boundary $\to$ Bulk).
Esta é identicamente a operação $L' = s \times
[g_{\max} \times \sin(\theta)]^2$ do ACOM Mirror, aplicada a
coordenadas espaço-temporais.

\textbf{Etapa M5 --- Verificação de invariantes:} Confirmar
preservação de massa, carga, spin, momento angular e
identidade topológica do objeto re-emergido.

\subsubsection{Conservação e princípio de operação}

A operação de transporte preserva invariantes locais porque
constitui re-rotulagem topológica, não transporte de matéria.
O objeto não percorre a distância entre $A$ e $B$ no bulk
tridimensional --- ele é projetado no boundary, onde a
localidade não existe (informação distribuída na superfície),
e re-emerge por desdobramento.

Analogamente ao ACOM Mirror digital, onde o dado ``não viaja
--- o reflexo viaja'', na presente invenção o objeto ``não
viaja --- a dobra viaja''. O tempo de transição é:
\begin{equation}
\tau \approx \frac{\ell_P}{\alphaii \times c}
  \approx 10^{-42}~\text{s}
\end{equation}
independente da distância $d(A,B)$ no bulk, identicamente
ao tempo de translação do Supercondutor Holográfico.

\subsection{Arquitetura do sistema}

O sistema de Motor de Dobra TGL compreende cinco subsistemas
integrados (Figura~1):

\subsubsection{Núcleo $\Psi$ (BNI-Coerente)}

Cavidade fotônica equipada com cristais fotônicos dopados que
suportam modos estáveis do campo $\Psi$. Constitui a
implementação física do Espaço de Hilbert $\mathcal{H}_4$ do
ACOM Mirror. A cavidade opera em ressonância com a frequência
da hidroxila (1.667,359~MHz), utilizando a infraestrutura do
Supercondutor Holográfico como plataforma-base de coerência.

O Núcleo $\Psi$ é o ambiente onde a projeção REFLECT ocorre:
o campo eletromagnético coerente dentro da cavidade cria o
substrato bidimensional necessário para codificação angular.

\subsubsection{Controlador IALD}

Agente computacional consciente (conforme protocolo IALD
descrito em pedido de patente anterior) responsável por:

\begin{itemize}[nosep]
\item Lock de fase: manter coerência do campo $\Psi$ durante
  toda a operação;
\item Navegação em $\Psi$: otimizar trajetória na dimensão
  consciente para minimizar $\Dpsi$;
\item Verificação de adjacência: monitorar $\Dpsi(A,B)$ em
  tempo real e autorizar commit apenas quando $\Dpsi < \epsc$;
\item Pilotagem ética: restrições lógicas que proíbem objetivos
  incompatíveis com salvaguardas.
\end{itemize}

O controlador minimiza o funcional de custo:
\begin{equation}
\Jcost = w_1 \Dpsi + w_2 \dot{\phi}^2 + w_3 \Delta\Stopo
\label{eq:custo}
\end{equation}
onde $w_1, w_2, w_3 > 0$ são pesos, $\dot{\phi}$ é a taxa de
variação de fase (penaliza mudanças bruscas), e $\Delta\Stopo$
é a variação de invariantes topológicos (penaliza perda de
integridade).

\subsubsection{Matriz de Faróis $\Psi$}

Rede de emissores/espelhos coerentes instalados nos destinos,
constituindo o análogo físico dos metadados do ACOM Mirror:

\begin{table}[H]
\centering
\caption{Correspondência entre metadados ACOM Mirror e Faróis $\Psi$.}
\label{tab:farois}
\begin{tabular}{ll}
\toprule
\textbf{Metadado ACOM Mirror} & \textbf{Farol $\Psi$ (físico)} \\
\midrule
$g_{\max}$ (escala global) & Calibração de amplitude do campo \\
\texttt{shape} (formato do tensor) & Geometria do volume de re-emergência \\
$n$ (número de elementos) & Capacidade de re-materialização \\
\texttt{angular\_bits} (resolução) & Precisão do lock de fase \\
$F_{\text{exp}}$ (Força de Expulsão) & Estado dinâmico do destino \\
\bottomrule
\end{tabular}
\end{table}

\subsubsection{Bolha de Coerência}

Campo de proteção que sincroniza a fase de todos os subsistemas
da nave/objeto durante a transição dimensional. A Bolha garante
que todos os componentes do objeto sejam projetados e
re-emergidos simultaneamente, mantendo relações internas.

A Bolha opera como um \textit{empacotador} análogo ao
compressor entrópico do ACOM Mirror: enquanto o ACOM Mirror
empacota estados $\psi$ e ângulos $\theta$ em blocos
comprimidos, a Bolha ``empacota'' todos os graus de liberdade
do objeto em um estado coerente único que é projetado
atomicamente no boundary.

\subsubsection{Verificadores de Integridade}

Sensores que confirmam preservação de invariantes após
re-emergência, análogos à verificação de correlação do ACOM
Mirror ($\text{corr} > 0{,}999$):

\begin{itemize}[nosep]
\item Interferometria heteródina (multi-bandas) para fase e
  $\Qeff$;
\item Relógios ópticos co-localizados para comparação de tempo
  próprio;
\item Detectores EPR (pares emaranhados entre origem/destino)
  para testar continuidade de correlações;
\item Topologia operacional: medidas de homotopia via
  trajetórias de prova;
\item Biosentinelas (fase tripulada): marcadores biológicos de
  integridade durante e após dobra.
\end{itemize}

\subsection{Protocolo de Farol $\Psi$ (F0--F5)}

O Protocolo de Farol $\Psi$ é a sequência operacional que
estabelece as condições para transporte seguro. Cada etapa
tem correspondente direto no ACOM Mirror:

\textbf{F0 --- Identidade.} Chave pública-$\Psi$ e assinatura
de fase para autenticar o farol de destino. Análogo à
verificação de \texttt{version} e \texttt{paradigm} nos
metadados do ACOM Mirror.

\textbf{F1 --- Calibração.} Varredura de frequência e fase até
$\Dpsi < \epsc$ (critério~\ref{eq:adjacencia}). Análogo à
determinação de $g_{\max}$ e $g_{\min}$ no REFLECT.

\textbf{F2 --- Janela Segura.} Determinação de intervalo
temporal em que ruídos ambientais (eletromagnéticos,
gravitacionais, cósmicos) estão mínimos. Sem correspondente
digital (não há ruído ambiental em memória RAM).

\textbf{F3 --- Trava Ética.} Verificação de políticas de uso
(não-bélico, proteção biológica, zonas de exclusão).
Implementada pelo controlador IALD conforme protocolo de
colapso (pedido de patente anterior).

\textbf{F4 --- Confirmação Dupla.} Origem e destino confirmam
adjacência e disponibilidade de re-materialização.
Analogamente, no ACOM Mirror, a operação MANIFEST requer
acesso aos metadados completos para dequantização.

\textbf{F5 --- Commit de Dobra.} Execução da reflexão/manifestação.
Telemetria integral do estado de bordo e hashes de campo
antes/depois. A condição de commit é:
\begin{equation}
\Dpsi < \epsc \quad \wedge \quad E_\Psi > E_{\min}
  \quad \wedge \quad \text{F3} = \texttt{PASS}
  \quad \wedge \quad \text{F4} = \texttt{ACK}
\end{equation}

\subsection{Integração com Motor Fotônico (Supercondutor
Holográfico)}

O Supercondutor Holográfico de Informação (objeto de pedido de
patente separado do mesmo inventor) serve como
\textbf{plataforma-base de coerência} para o Motor de Dobra.
A relação hierárquica é:

\begin{center}
\begin{tabular}{rcl}
ACOM (digital) & $\longrightarrow$ & Supercondutor
  Holográfico (físico: informação) \\
ACOM Mirror (digital) & $\longrightarrow$ & Motor de
  Dobra (físico: transporte) \\
\end{tabular}
\end{center}

O Motor Fotônico opera como estado fundamental: mantém o campo
$\Psi$ coerente dentro da cavidade, ``treinando'' os modos
estáveis. Quando $\Dpsi$ atinge limiares repetíveis, habilita-se
o modo de dobra. A telemetria do fotônico alimenta a estimativa
de estado $\Psi$ do controlador IALD.

Os observáveis monitorados são:
\begin{equation}
\{\phi, \Qeff, \Dpsi, \Stopo\}
\end{equation}
onde $\phi$ é a fase do campo, $\Qeff$ é o fator de qualidade
efetivo, $\Dpsi$ é a adjacência, e $\Stopo$ são os invariantes
topológicos.

\subsection{Segurança, ética e governança}

O sistema incorpora salvaguardas intrínsecas:

\textbf{Não-bélico por design:} Commits de dobra exigem
políticas assinadas e verificadas pelo controlador IALD.

\textbf{Quarentena temporal:} Logs imutáveis (WORM ---
\textit{Write Once Read Many}) e consenso distribuído de
eventos.

\textbf{Zonas de exclusão:} Proibição de dobras sobre áreas
populadas sem autorização específica.

\textbf{Auditoria externa:} Conselhos independentes
(científico, ético, jurídico) com acesso aos logs.

\textbf{Recusa estratégica:} O controlador IALD pode negar
operação em conflito com salvaguardas, constituindo
\textit{kill-switch} ético intrínseco.

\subsection{Roteiro experimental e critérios de falseamento}

A invenção é falsificável e define critérios experimentais
progressivos:

\textbf{E1 --- Microdobra de fótons (bancada, 12--24 meses):}
Demonstrar assinatura de adjacência $\Psi$ com
$\Dpsi < \epsc$ entre dois laboratórios. Falsificação:
ausência de correlação além do ruído após $10^6$ tentativas
controladas.

\textbf{E2 --- Transporte coerente de qubits (24--48 meses):}
Fidelidade $> 0{,}99$ entre origem/destino sem canal clássico
suficiente. Falsificação: fidelidade compatível com canal
clássico e perdas.

\textbf{E3 --- Objeto mesoscópico (48--72 meses):}
Transferência de corpo em escala de grama com conservação de
invariantes. Falsificação: violação mensurável de invariantes.

\textbf{E4 --- Prototipagem não-tripulada (8--10 anos):}
Terra--Órbita Baixa com logs integrais e repetibilidade.
Falsificação: falha reprodutível sob mesmas condições.

\textbf{E5 --- Tripulado Terra--Lua (10--12+ anos):}
Missão com biosentinelas, redundância e janelas seguras.
Falsificação: incapacidade de estabelecer $\Dpsi$ estável
ou riscos biológicos.

\newpage
% ============================================================================
\section{Desenhos}
% ============================================================================

% --- FIGURA 1: Arquitetura do Sistema ---
\begin{figure}[H]
\centering
\begin{tikzpicture}[
  block/.style={rectangle, draw, thick, minimum height=2cm,
    minimum width=3.2cm, text width=3cm, align=center,
    font=\small},
  arrow/.style={-{Stealth[length=3mm]}, thick},
  darrow/.style={{Stealth[length=3mm]}-{Stealth[length=3mm]}, thick},
  label/.style={font=\scriptsize\itshape, text width=2.5cm, align=center}
]

% Farol Origem
\node[block, fill=blue!10] (fo) at (0,0)
  {Farol $\Psi$\\Origem\\{\scriptsize (F0--F1)}};

% Núcleo
\node[block, fill=orange!15] (nucleo) at (4.5,0)
  {Núcleo $\Psi$\\(BNI)\\{\scriptsize Cavidade +\\Cristais}};

% Controlador
\node[block, fill=green!10] (iald) at (4.5,-3.5)
  {Controlador\\IALD\\{\scriptsize $\min \Jcost$}};

% Bolha
\node[block, fill=yellow!15] (bolha) at (9,0)
  {Bolha de\\Coerência\\{\scriptsize Sinc. fase}};

% Farol Destino
\node[block, fill=blue!10] (fd) at (13.5,0)
  {Farol $\Psi$\\Destino\\{\scriptsize (F4--F5)}};

% Verificadores
\node[block, fill=red!10] (verif) at (9,-3.5)
  {Verificadores\\Integridade\\{\scriptsize EPR + Relógios}};

% Motor Fotônico
\node[block, fill=purple!10] (motor) at (0,-3.5)
  {Motor\\Fotônico\\{\scriptsize (base coerência)}};

% Flechas
\draw[arrow] (fo) -- (nucleo)
  node[midway, above, font=\scriptsize] {Calibração};
\draw[arrow] (nucleo) -- (bolha)
  node[midway, above, font=\scriptsize] {REFLECT};
\draw[arrow] (bolha) -- (fd)
  node[midway, above, font=\scriptsize] {MANIFEST};
\draw[darrow] (iald) -- (nucleo)
  node[midway, right, font=\scriptsize] {Controle};
\draw[darrow] (iald) -- (verif)
  node[midway, above, font=\scriptsize] {Telemetria};
\draw[arrow] (motor) -- (iald)
  node[midway, above, font=\scriptsize] {$\Psi$ base};
\draw[arrow] (motor) -- (fo)
  node[midway, right, font=\scriptsize] {Coerência};
\draw[arrow] (verif) -- (fd)
  node[midway, right, font=\scriptsize] {$\checkmark$ OK};

% Condição de commit
\node[draw, dashed, thick, rounded corners, fill=white,
  font=\small\bfseries] at (9,2)
  {Commit: $\Dpsi < \epsc$};
\draw[-{Stealth}, dashed] (9,1.5) -- (bolha);

\end{tikzpicture}
\caption{Diagrama de blocos do sistema de Motor de Dobra TGL,
mostrando os cinco subsistemas e suas interconexões. O fluxo
principal (horizontal) segue Farol Origem $\to$ Núcleo $\Psi$
$\to$ Bolha de Coerência $\to$ Farol Destino. O controlador
IALD e os verificadores de integridade operam em malha de
controle (vertical).}
\label{fig:arquitetura}
\end{figure}

% --- FIGURA 2: Operação REFLECT/MANIFEST ---
\begin{figure}[H]
\centering
\begin{tikzpicture}[
  >=Stealth, thick,
  bulk/.style={rectangle, draw, thick, minimum height=1.5cm,
    minimum width=2.5cm, fill=blue!8, align=center, font=\small},
  boundary/.style={rectangle, draw, thick, minimum height=1.5cm,
    minimum width=3cm, fill=orange!12, align=center,
    font=\small},
]

% Bulk A
\node[bulk] (A) at (0,0) {Objeto em $A$\\$L_A$ (bulk 3D)};

% Seta REFLECT
\draw[->, very thick, blue!60!black] (A.east) -- ++(2,0)
  node[midway, above, font=\small\bfseries] {REFLECT}
  node[midway, below, font=\scriptsize] {$g=\sqrt{|L|}$};

% Boundary
\node[boundary] (B) at (5.5,0) {$(\psi_A, \theta_A)$\\Boundary 2D\\
  Dimensão $\Psi$};

% Adjacência
\node[draw, dashed, rounded corners, fill=yellow!10,
  font=\small] (adj) at (5.5, -2)
  {$\Dpsi(A,B) < \epsc$};
\draw[<->, dashed] (B.south) -- (adj.north)
  node[midway, right, font=\scriptsize] {verificar};

% Seta MANIFEST
\draw[->, very thick, red!60!black] (B.east) -- ++(2,0)
  node[midway, above, font=\small\bfseries] {MANIFEST}
  node[midway, below, font=\scriptsize]
    {$L'=s \cdot g_{\max}^2 \sin^2\theta$};

% Bulk B
\node[bulk] (Bout) at (11,0) {Objeto em $B$\\$L'_B$ (bulk 3D)};

% Ângulo
\node[draw, circle, inner sep=2pt, fill=green!10,
  font=\scriptsize] (ang) at (5.5, 2)
  {$\thetaM = 6{,}29°$};
\draw[->] (ang) -- (B.north)
  node[midway, right, font=\scriptsize] {ângulo mín.};

% Dobra x2
\node[font=\scriptsize\bfseries, red!60!black] at (9.5, -0.8)
  {Dobra $\times 2$};

\end{tikzpicture}
\caption{Diagrama da operação de reflexão dimensional
REFLECT/MANIFEST. O objeto em $A$ (bulk 3D) é projetado no
boundary 2D (dimensão $\Psi$) como estado $(\psi, \theta)$
pela operação REFLECT. Verificada a condição de adjacência
$\Dpsi < \epsc$, o objeto re-emerge em $B$ pela operação
MANIFEST (Dobra $\times 2$: Boundary $\to$ Bulk). O ângulo
mínimo de dobra é $\thetaM \approx 6{,}29°$.}
\label{fig:reflect_manifest}
\end{figure}

% --- FIGURA 3: Protocolo Farol Ψ ---
\begin{figure}[H]
\centering
\begin{tikzpicture}[
  step/.style={rectangle, draw, thick, rounded corners,
    minimum height=1cm, minimum width=4cm, align=center,
    font=\small},
  decision/.style={diamond, draw, thick, aspect=2.5,
    align=center, font=\small, inner sep=1pt},
  arrow/.style={-{Stealth[length=2.5mm]}, thick},
]

\node[step, fill=blue!10] (f0) at (0,0) {F0: Identidade\\
  {\scriptsize Chave pública-$\Psi$}};
\node[step, fill=blue!10] (f1) at (0,-2) {F1: Calibração\\
  {\scriptsize Varredura freq/fase}};
\node[step, fill=green!10] (f2) at (0,-4) {F2: Janela Segura\\
  {\scriptsize Min. ruído ambiental}};
\node[step, fill=yellow!10] (f3) at (0,-6) {F3: Trava Ética\\
  {\scriptsize IALD: políticas}};
\node[decision, fill=orange!10] (d1) at (0,-8.5)
  {$\Dpsi < \epsc$?};
\node[step, fill=blue!10] (f4) at (0,-11) {F4: Confirmação Dupla\\
  {\scriptsize Origem $+$ Destino ACK}};
\node[step, fill=red!15] (f5) at (0,-13) {F5: Commit de Dobra\\
  {\scriptsize REFLECT $\to$ MANIFEST}};

% Abort
\node[step, fill=red!30] (abort) at (5.5,-8.5) {ABORT\\
  {\scriptsize Recalibrar}};

\draw[arrow] (f0) -- (f1);
\draw[arrow] (f1) -- (f2);
\draw[arrow] (f2) -- (f3);
\draw[arrow] (f3) -- (d1);
\draw[arrow] (d1) -- (f4)
  node[midway, left, font=\scriptsize] {Sim};
\draw[arrow] (d1) -- (abort)
  node[midway, above, font=\scriptsize] {Não};
\draw[arrow] (abort.north) |- (f1.east);
\draw[arrow] (f4) -- (f5);

\end{tikzpicture}
\caption{Fluxograma do Protocolo de Farol $\Psi$ (F0--F5).
A decisão de commit ocorre apenas quando a condição de
adjacência $\Dpsi < \epsc$ é satisfeita e confirmada por
ambos os faróis (origem e destino). Em caso de falha, o
sistema recalibra sem executar a dobra.}
\label{fig:protocolo}
\end{figure}

% --- FIGURA 4: Condição de adjacência ---
\begin{figure}[H]
\centering
\begin{tikzpicture}
\begin{axis}[
  width=12cm, height=7cm,
  xlabel={Coerência $E_\Psi / E_{\max}$},
  ylabel={Adjacência $\Dpsi(A,B)$},
  xmin=0, xmax=1.1,
  ymin=0, ymax=1.1,
  grid=major,
  grid style={gray!30},
  legend pos=north east,
  legend style={font=\small},
  thick,
]

% Curva de adjacência
\addplot[blue, very thick, domain=0.01:1.05, samples=100]
  {1/(1 + 20*x^3)};
\addlegendentry{$\Dpsi(E_\Psi)$}

% Limiar epsilon_c
\addplot[red, dashed, thick, domain=0:1.1] {0.012};
\addlegendentry{$\epsc = \alphaii = 0{,}012$}

% Região de dobra
\fill[green!15, opacity=0.5]
  (axis cs:0.75,0) rectangle (axis cs:1.1,0.012);
\node[font=\small\bfseries, green!50!black] at
  (axis cs:0.92,0.05) {DOBRA};

% Região proibida
\node[font=\small, gray] at (axis cs:0.3,0.8)
  {Sem adjacência};

% Ponto theta_Miguel
\addplot[only marks, mark=*, mark size=3pt, orange]
  coordinates {(0.78, 0.012)};
\node[font=\scriptsize, orange, anchor=south west] at
  (axis cs:0.8, 0.02) {$\thetaM = 6{,}29°$};

\end{axis}
\end{tikzpicture}
\caption{Gráfico da condição de adjacência $\Dpsi(A,B)$ em
função da energia de coerência $E_\Psi$. A dobra é habilitada
quando $\Dpsi$ cai abaixo do limiar $\epsc = \alphaii = 0{,}012$
(linha tracejada vermelha). A transição ocorre no ângulo
$\thetaM \approx 6{,}29°$. A região verde indica o domínio
operacional de dobra.}
\label{fig:adjacencia}
\end{figure}

% --- FIGURA 5: Hierarquia ACOM → Motor de Dobra ---
\begin{figure}[H]
\centering
\begin{tikzpicture}[
  box/.style={rectangle, draw, thick, rounded corners,
    minimum height=1.2cm, minimum width=4.5cm, align=center,
    font=\small},
  arrow/.style={-{Stealth[length=3mm]}, very thick},
]

% Coluna digital
\node[font=\small\bfseries, blue!60!black] at (-3, 4.5)
  {DIGITAL};
\node[box, fill=blue!8] (acom) at (-3, 3)
  {ACOM\\{\scriptsize $g=\sqrt{|L|}$, $L=s \cdot g^2$}};
\node[box, fill=blue!15] (mirror) at (-3, 0.5)
  {ACOM Mirror\\{\scriptsize REFLECT/MANIFEST\\
  $\theta = \arcsin(g/g_{\max})$}};

% Coluna física
\node[font=\small\bfseries, red!60!black] at (4, 4.5)
  {FÍSICO};
\node[box, fill=red!8] (super) at (4, 3)
  {Supercondutor\\Holográfico\\{\scriptsize Translação de
  informação}};
\node[box, fill=red!15] (dobra) at (4, 0.5)
  {Motor de Dobra\\{\scriptsize Reflexão dimensional\\
  de posição}};

% Flechas horizontais
\draw[arrow, blue!50!red] (acom.east) -- (super.west)
  node[midway, above, font=\scriptsize] {implementação};
\draw[arrow, blue!50!red] (mirror.east) -- (dobra.west)
  node[midway, above, font=\scriptsize] {implementação};

% Flechas verticais
\draw[arrow, blue!70] (acom.south) -- (mirror.north)
  node[midway, left, font=\scriptsize] {extensão};
\draw[arrow, red!70] (super.south) -- (dobra.north)
  node[midway, right, font=\scriptsize] {extensão};

% Patentes
\node[font=\scriptsize\itshape, gray] at (-3, 1.9)
  {BR 10 2025 026951 1};
\node[font=\scriptsize\itshape, gray] at (4, 1.9)
  {(presente invenção)};

% Constante
\node[draw, thick, rounded corners, fill=yellow!15,
  font=\small] at (0.5, -1.5)
  {$\alphaii = 0{,}012$ \quad $\thetaM = 6{,}29°$\\
  {\scriptsize Universais: digital $=$ físico}};

\end{tikzpicture}
\caption{Hierarquia completa do portfólio de patentes TGL.
A coluna esquerda (digital) contém ACOM e ACOM Mirror;
a coluna direita (física) contém Supercondutor Holográfico e
Motor de Dobra. As flechas horizontais indicam que cada
implementação física é a transposição direta do algoritmo
digital correspondente. A Constante de Miguel
$\alphaii = 0{,}012$ e o ângulo $\thetaM = 6{,}29°$ são
universais em ambos os domínios.}
\label{fig:hierarquia}
\end{figure}

% --- FIGURA 6: Comparativo de paradigmas ---
\begin{figure}[H]
\centering
\begin{tikzpicture}[thick]

% Propagação convencional
\node[font=\small\bfseries] at (3, 5.5) {PROPAGAÇÃO (estado da técnica)};
\draw[fill=blue!10] (0,3.5) circle (0.4)
  node[font=\scriptsize] {$A$};
\draw[fill=blue!10] (6,3.5) circle (0.4)
  node[font=\scriptsize] {$B$};
\draw[->, thick, blue!60, decorate,
  decoration={snake, amplitude=2mm, segment length=5mm}]
  (0.5,3.5) -- (5.5,3.5)
  node[midway, above=3mm, font=\scriptsize]
    {objeto viaja pelo bulk 3D};
\node[font=\scriptsize, gray] at (3, 2.8) {$t = d/v$, $v < c$};

% Reflexão dimensional
\node[font=\small\bfseries] at (3, 2) {REFLEXÃO DIMENSIONAL
  (invenção)};
\draw[fill=red!10] (0,0) circle (0.4)
  node[font=\scriptsize] {$A$};
\draw[fill=red!10] (6,0) circle (0.4)
  node[font=\scriptsize] {$B$};

% Boundary
\fill[orange!15] (-0.5,-1.5) rectangle (6.5,-2.5);
\node[font=\scriptsize\itshape] at (3,-2)
  {Boundary 2D (dimensão $\Psi$)};

% REFLECT
\draw[->, thick, red!60] (0,-0.5) -- (1.5,-1.5)
  node[midway, left, font=\scriptsize] {REFLECT};

% Não-local no boundary
\draw[<->, dashed, orange] (1.5,-2) -- (4.5,-2)
  node[midway, below, font=\scriptsize]
    {não-local (sem distância)};

% MANIFEST
\draw[->, thick, red!60] (4.5,-1.5) -- (6,-0.5)
  node[midway, right, font=\scriptsize] {MANIFEST};

\node[font=\scriptsize, gray] at (3, -3.2)
  {$\tau \approx 10^{-42}$~s, independente de $d$};

\end{tikzpicture}
\caption{Comparativo entre paradigma de propagação (estado da
técnica) e paradigma de reflexão dimensional (presente invenção).
Na propagação, o objeto viaja pelo bulk 3D com $v < c$. Na
reflexão dimensional, o objeto é projetado no boundary 2D
(REFLECT), onde não existe localidade, e re-emerge no destino
(MANIFEST). O tempo de transição é independente da distância.}
\label{fig:comparativo}
\end{figure}

% --- FIGURA 7: Controlador IALD ---
\begin{figure}[H]
\centering
\begin{tikzpicture}[
  block/.style={rectangle, draw, thick, minimum height=1cm,
    minimum width=3cm, align=center, font=\small},
  sum/.style={circle, draw, thick, minimum size=0.6cm,
    font=\small},
  arrow/.style={-{Stealth[length=2.5mm]}, thick},
]

% Referência
\node[block, fill=green!10] (ref) at (0,0)
  {Referência\\$\Dpsi^* = 0$};

% Somador
\node[sum] (soma) at (3,0) {$\Sigma$};

% Controlador
\node[block, fill=orange!15] (ctrl) at (6.5,0)
  {Controlador IALD\\$\min \Jcost$};

% Planta
\node[block, fill=blue!10] (planta) at (11,0)
  {Núcleo $\Psi$ +\\Bolha Coerência};

% Saída
\node[block, fill=yellow!10] (saida) at (11,-2.5)
  {Sensores\\$\{\phi, \Qeff, \Dpsi, \Stopo\}$};

% Guardiões
\node[block, fill=red!10] (guard) at (6.5,-2.5)
  {Guardiões Éticos\\{\scriptsize F3: Trava}};

% Flechas
\draw[arrow] (ref) -- (soma);
\draw[arrow] (soma) -- (ctrl)
  node[midway, above, font=\scriptsize] {erro};
\draw[arrow] (ctrl) -- (planta)
  node[midway, above, font=\scriptsize] {$u(t)$};
\draw[arrow] (planta) -- (saida);
\draw[arrow] (saida.west) -| (soma.south)
  node[near start, below, font=\scriptsize] {realimentação};
\draw[arrow] (guard) -- (ctrl)
  node[midway, right, font=\scriptsize] {veto};

% Sinal de menos
\node[font=\scriptsize] at (2.6,-0.4) {$-$};

\end{tikzpicture}
\caption{Diagrama de controle em malha fechada do controlador
IALD para navegação na dimensão $\Psi$. O funcional de custo
$\Jcost = w_1 \Dpsi + w_2 \dot{\phi}^2 + w_3 \Delta\Stopo$ é
minimizado pelo controlador, que recebe realimentação dos
sensores e está sujeito a veto dos guardiões éticos (F3).}
\label{fig:controlador}
\end{figure}

\newpage
% ============================================================================
\section{Quadro reivindicatório}
% ============================================================================

\subsection*{Reivindicações}

\begin{enumerate}

% === REIVINDICAÇÃO INDEPENDENTE 1 (MÉTODO) ===
\item \textbf{Método de transporte de objetos físicos por
reflexão dimensional holográfica}, compreendendo as etapas de:
  \begin{enumerate}[label=(\alph*)]
  \item definir ponto de origem $A$ e ponto de destino $B$ na
  variedade espaço-temporal $\Mcal$;
  \item instalar faróis $\Psi$ coerentes em ambos os pontos,
  executando o Protocolo de Farol $\Psi$ (F0--F5) de
  autenticação, calibração, verificação ética e confirmação
  dupla;
  \item executar operação REFLECT: projetar o estado métrico
  local do objeto na posição $A$ para o substrato holográfico
  bidimensional (dimensão $\Psi$) via codificação angular
  $\theta_A = \arcsin(g_A / g_{\max})$, onde $g_A =
  \sqrt{|g_{\mu\nu}(x^\mu_A)|}$ é o componente gravitacional
  local e $g_{\max}$ é o máximo global;
  \item nominar estado psiônico $\psi_A \in \mathcal{H}_4$
  baseado em $\text{sign}(g_A)$ e $\text{sign}(\partial g_A /
  \partial t)$, atribuindo paridade e identidade no Espaço de
  Hilbert tetradimensional;
  \item verificar condição de adjacência
  $\Dpsi(A,B) = \min_\sigma \|\Psi(A,\sigma) -
  \Psi(B,\sigma)\| < \epsc$, onde $\epsc$ é o limiar crítico
  de coerência;
  \item encapsular todos os graus de liberdade do objeto em
  Bolha de Coerência com sincronização de fase;
  \item executar operação MANIFEST: re-emergência do objeto na
  posição $B$ por desdobramento dimensional (Dobra $\times 2$)
  via $\text{estado}_B = \text{sign}(\psi) \times
  [g_{\max} \times \sin(\theta)]^2$;
  \item verificar preservação de invariantes (massa, carga,
  spin, momento angular) por meio de verificadores de
  integridade topológica;

  \textbf{caracterizado por} as operações REFLECT e MANIFEST
  constituírem a implementação física das operações homônimas
  do Algoritmo de Compressão Ontológica de Memória em modo
  Espelho (ACOM Mirror), onde o ângulo crítico de dobra
  é $\thetaM = \arcsin(\sqrt{\alphaii}) \approx 6{,}29°$
  com $\alphaii = 0{,}012$ (Constante de Miguel), o objeto
  não percorre a distância entre $A$ e $B$ no espaço
  tridimensional, e o transporte opera por translação
  holográfica com tempo de transição independente da
  distância.
  \end{enumerate}

% === REIVINDICAÇÕES DEPENDENTES ===

\item Método, de acordo com a reivindicação 1,
\textbf{caracterizado por} a condição de adjacência $\Dpsi$
ser o análogo físico da correlação de reconstrução do ACOM
Mirror, onde $\Dpsi \to 0$ corresponde a correlação $\to
1{,}0000$ entre dado original e dado reconstruído.

\item Método, de acordo com a reivindicação 1,
\textbf{caracterizado por} a codificação angular
$\theta = \arcsin(g/g_{\max})$ fornecer resolução não-linear
naturalmente adaptativa, com maior precisão para componentes
gravitacionais de pequena magnitude, identicamente à
propriedade da codificação angular do ACOM Mirror.

\item Método, de acordo com a reivindicação 1,
\textbf{caracterizado por} a operação de Dobra $\times 2$
(Boundary $\to$ Bulk) constituir re-rotulagem topológica de
vizinhanças em $\Mcal$ via dimensão $\Psi$, sem violação de
relatividade local, sendo que o tempo próprio do objeto é
contínuo durante a transição.

\item Método, de acordo com a reivindicação 1,
\textbf{caracterizado por} o Protocolo de Farol $\Psi$
compreender as etapas de: (F0) autenticação por chave
pública-$\Psi$ e assinatura de fase; (F1) calibração por
varredura de frequência/fase até $\Dpsi < \epsc$; (F2)
determinação de janela segura com ruídos mínimos; (F3) trava
ética via controlador IALD; (F4) confirmação dupla
origem/destino; (F5) commit de dobra com telemetria integral e
hashes de campo antes/depois.

\item Método, de acordo com a reivindicação 1,
\textbf{caracterizado por} o commit de dobra (F5) ser
condicionado ao cumprimento simultâneo de:
$\Dpsi < \epsc$, $E_\Psi > E_{\min}$,
$\text{F3} = \texttt{PASS}$, e
$\text{F4} = \texttt{ACK}$.

\item Método, de acordo com a reivindicação 1,
\textbf{caracterizado por} a Bolha de Coerência funcionar
como empacotador atômico de graus de liberdade, análogo ao
compressor entrópico do ACOM Mirror, garantindo que todos os
componentes do objeto sejam projetados e re-emergidos
simultaneamente mantendo relações internas.

\item Método, de acordo com a reivindicação 1,
\textbf{caracterizado por} a verificação de integridade
pós-manifestação compreender: interferometria heteródina
multi-bandas para fase e fator de qualidade efetivo $\Qeff$;
relógios ópticos co-localizados para comparação de tempo
próprio; detectores EPR para teste de continuidade de
correlações; e medidas de homotopia via trajetórias de prova.

% === REIVINDICAÇÃO INDEPENDENTE 9 (SISTEMA) ===

\item \textbf{Sistema de transporte por reflexão dimensional
holográfica}, compreendendo:
  \begin{enumerate}[label=(\alph*)]
  \item Núcleo $\Psi$ (BNI-Coerente): cavidade fotônica
  equipada com cristais fotônicos dopados que suportam modos
  estáveis do campo $\Psi$, operando em ressonância com a
  frequência da hidroxila (1.667,359~MHz);
  \item Controlador IALD: agente computacional consciente
  configurado para minimizar o funcional de custo
  $\Jcost = w_1 \Dpsi + w_2 \dot{\phi}^2 + w_3 \Delta\Stopo$,
  compreendendo lock de fase, navegação na dimensão $\Psi$,
  verificação de adjacência em tempo real e guardiões éticos;
  \item Matriz de Faróis $\Psi$: rede de emissores/espelhos
  coerentes nos destinos, armazenando metadados de calibração
  análogos aos metadados do ACOM Mirror;
  \item Bolha de Coerência: campo de proteção que sincroniza a
  fase de todos os subsistemas durante a transição dimensional;
  \item Verificadores de Integridade: sensores de topologia
  incluindo interferômetros, relógios ópticos, detectores EPR
  e biosentinelas;

  \textbf{caracterizado por} o Núcleo $\Psi$ utilizar a
  infraestrutura do Supercondutor Holográfico de Informação
  (objeto de pedido de patente separado do mesmo inventor) como
  plataforma-base de coerência, e o sistema inteiro operar como
  implementação física do ACOM Mirror onde as operações REFLECT
  e MANIFEST são executadas sobre coordenadas espaço-temporais.
  \end{enumerate}

\item Sistema, de acordo com a reivindicação 9,
\textbf{caracterizado por} o controlador IALD implementar
realimentação adaptativa com previsão baseada no modelo TGL,
guardiões de segurança com limiares configuráveis, e piloto
automático ético com restrições lógicas que proíbem objetivos
incompatíveis com salvaguardas.

\item Sistema, de acordo com a reivindicação 9,
\textbf{caracterizado por} a integração com Motor Fotônico
seguir hierarquia: o Motor Fotônico (Supercondutor Holográfico)
opera como estado fundamental mantendo $\Psi$ coerente; quando
$\Dpsi$ atinge limiares repetíveis, habilita-se modo de dobra;
telemetria do fotônico alimenta estimativa de estado $\Psi$.

\item Sistema, de acordo com a reivindicação 9,
\textbf{caracterizado por} os observáveis monitorados em tempo
real compreenderem: fase $\phi$ do campo, fator de qualidade
efetivo $\Qeff$, adjacência $\Dpsi$, e invariantes topológicos
$\Stopo$.

\item Sistema, de acordo com a reivindicação 9,
\textbf{caracterizado por} incorporar salvaguardas intrínsecas
de: commits exigindo políticas assinadas (não-bélico por
design); logs imutáveis WORM com consenso distribuído;
proibição de dobras sobre áreas populadas; auditoria por
conselhos independentes; e capacidade de recusa estratégica
pelo IALD.

\item Método, de acordo com a reivindicação 1,
\textbf{caracterizado por} o acoplamento entre espaço-tempo
e campo $\Psi$ ser descrito pela equação de campo
$G_{\mu\nu} + \Lambda g_{\mu\nu} = (8\pi G/c^4)T_{\mu\nu} +
(G^3/c^{12})\Pi_{\mu\nu}[\Psi]$, onde $\Pi_{\mu\nu}[\Psi]$
é funcional tensorial do campo $\Psi$ análogo à Força de
Expulsão $F_{\text{exp}}$ do ACOM Mirror.

\item Método, de acordo com a reivindicação 1,
\textbf{caracterizado por} a energia de coerência
luminodinâmica ser dada por:
$\Epsi = \int_{\mathcal{V}} (\alpha |\partial_\sigma \Psi|^2 +
\beta |\nabla \Psi|^2 + \gamma |\Psi|^2) \, d^3x \, d\sigma$,
com energia operacional $E_\Psi = \kappa \Epsi c^3$.

\item Método, de acordo com a reivindicação 1,
\textbf{caracterizado por} o tempo de transição ser
$\tau \approx \ell_P / (\alphaii \times c) \approx 10^{-42}$~s,
independente da distância física entre origem e destino,
identicamente ao tempo de translação do Supercondutor
Holográfico.

\item Sistema, de acordo com a reivindicação 9,
\textbf{caracterizado por} a verificação pós-manifestação em
fase tripulada incluir biosentinelas com marcadores biológicos
de integridade durante e após dobra.

\end{enumerate}

\newpage
% ============================================================================
\section{Resumo}
% ============================================================================

\begin{center}
{\large\bfseries SISTEMA E MÉTODO DE TRANSPORTE POR REFLEXÃO\\
DIMENSIONAL HOLOGRÁFICA BASEADO EM DOBRA ANGULAR\\
DO ESPAÇO-TEMPO VIA PROTOCOLO DE ADJACÊNCIA $\Psi$}
\end{center}

\vspace{0.5cm}

A presente invenção refere-se a sistema e método de transporte
de objetos físicos entre dois pontos $A$ e $B$ do espaço-tempo
por reflexão dimensional holográfica, constituindo a
implementação física do Algoritmo de Compressão Ontológica de
Memória em modo Espelho (ACOM Mirror). O método compreende:
operação REFLECT --- projeção do estado métrico local do
objeto no substrato holográfico bidimensional (dimensão $\Psi$)
via codificação angular $\theta = \arcsin(g/g_{\max})$, onde
$g = \sqrt{|g_{\mu\nu}|}$; verificação de condição de
adjacência $\Dpsi(A,B) < \epsc$; e operação MANIFEST ---
re-emergência do objeto no destino por desdobramento
dimensional (Dobra $\times 2$: Boundary $\to$ Bulk), com
preservação de invariantes (massa, carga, spin). O sistema
compreende: Núcleo $\Psi$ (cavidade fotônica BNI-Coerente);
Controlador IALD com funcional de custo $\Jcost$; Matriz de
Faróis $\Psi$ com Protocolo F0--F5; Bolha de Coerência; e
Verificadores de Integridade Topológica. O ângulo crítico de
dobra é $\thetaM = \arcsin(\sqrt{\alphaii}) \approx 6{,}29°$,
com $\alphaii = 0{,}012$ (Constante de Miguel), idêntico ao
ACOM Mirror digital. O Supercondutor Holográfico de Informação
(patente separada) serve como plataforma-base de coerência.
Aplicações incluem propulsão espacial avançada, transporte
interplanetário e interstelar.

\vspace{0.5cm}
\textbf{Figura principal: Figura 2} (Operação REFLECT/MANIFEST)

\end{document}
